\section{AutoEncoder 候補}

\subsection{再構成誤差による検出}

長さ $N$ の入力列 $X_k$ に対して、AE は
$\hat{X}_k=f_{\theta}(X_k)$ を返す。anomaly score は、特徴群 $g$ ごとの
重み付き再構成誤差として式 \eqref{eq:score} で定義する。

\begin{equation}
  s(X_k) = \sum_{g \in G} w_g
         \frac{1}{N |I_g|}
         \sum_{i=1}^{N}
         \sum_{j \in I_g}
         \ell \left(X_{k,i,j}, \hat{X}_{k,i,j}\right).
  \label{eq:score}
\end{equation}

初期実装ではすべての入力が scalar なので、loss は MSE または Huber loss に固定する。
histogram profile を将来追加する場合だけ、Jensen-Shannon divergence などの
分布用 loss を別 ablation として導入する。しきい値 $\gamma$ は
benign calibration split の上位分位点から決め、評価時には
false alarms per volume per day を主指標にする。

\subsection{候補アーキテクチャ}

ユーザ仮説では、CNN が I/O 情報の相関や局所特徴を抽出し、GRU が時系列文脈を
付与し、decoder は bottleneck 表現から入力列を再構成する。これは自然な候補であるが、
本稿で AE を選ぶ理由は、複雑な深層モデルを使いたいからではない。
教師なしにより未知脅威を正常時系列からの逸脱として扱えること、かつ
固定 shape の再構成問題として MNN に載せやすいことが主理由である。
volume あたり 500 KB の detector-data 制約、さらに多数 volume へ同じ検出器を展開する制約では、
小型候補との比較が必須である。
本稿では operator の役割を固定する。Dense は情報圧縮、GRU は各 frame への
時系列文脈付与、Conv1D は scalar 入力を複数の局所時間視点へ展開して短距離
pattern を抽出する処理とする。したがって、GRU や Conv1D に圧縮責任を曖昧に
負わせず、圧縮は明示的な Dense bottleneck に置く。

\begin{table}[tb]
  \centering
  \caption{評価すべき AE 候補}
  \label{tab:model-candidates}
  \begin{tabularx}{\linewidth}{lY Y}
    \toprule
    候補 & 利点 & 主なリスク \\
    \midrule
    MLP bottleneck AE &
    最小の operator set で MNN 変換が容易。メモリ見積もりも単純。 &
    時系列構造を平坦化するため、ゆっくりした攻撃や順序依存を失う可能性。 \\
    GRU contextual AE &
    GRU で各 frame に temporal context を付与し、Dense bottleneck で明示的に圧縮できる。 &
    recurrent operator の workspace と量子化 drift を要測定。 \\
    Temporal convolution AE &
    Conv1D で局所時間 pattern を複数チャネルへ展開し、Dense bottleneck で圧縮するため、
    計算量と activation memory を制御しやすい。 &
    長い文脈や不規則な attack rhythm への追従が弱い可能性。 \\
    Tiny CNN-GRU AE &
    10 秒フレーム列の局所変化を Conv1D で展開・抽出し、GRU で文脈化し、Dense bottleneck で圧縮できる。 &
    最も複雑で、volume あたり 500 KB detector-data と MNN operator support の制約を受けやすい。 \\
    \bottomrule
  \end{tabularx}
\end{table}

表 \ref{tab:concrete-models} に、$D=12$、$N=12$、すなわち 144 次元の
flattened 入力を仮定した具体的な小型 AE 例を示す。これは重みだけの粗い
見積もりであり、入力統計/state と transient scratch は別途測定する必要がある。

\begin{table}[tb]
  \centering
  \caption{メモリ制約を考慮した具体 AE 例}
  \label{tab:concrete-models}
  \small
  \resizebox{\linewidth}{!}{%
  \begin{tabular}{llrrl}
    \toprule
    ID & Model sketch & Params & FP32 & Role \\
    \midrule
    AE-0 & \texttt{144-32-8-32-144} & 9,944 & 39 KB & flat baseline \\
    AE-1 & \texttt{144-64-16-64-144} & 20,768 & 83 KB & flat capacity check \\
    AE-2 & two-level dense bottleneck & 5,636 & 23 KB & explicit dense compression \\
    AE-3a & GRU context + dense bottleneck, hidden 24 & 7,052 & 29 KB & recurrent context AE \\
    AE-3b & GRU context + dense bottleneck, hidden 32 & 11,700 & 47 KB & stronger recurrent AE \\
    AE-4 & temporal Conv1D AE, 24 channels & 5,684 & 23 KB & MNN-friendly temporal AE \\
    AE-5 & Conv1D + GRU context + dense bottleneck & 8,804 & 36 KB & constrained CNN-GRU test \\
    \bottomrule
  \end{tabular}
  }
\end{table}

初期実装では AE-0、AE-2、AE-4、AE-3a の順で評価する。AE-0 は性能を期待する
本命ではなく、教師なし再構成がこの観測境界で最低限の信号を持つかを見る
flattened AE の最低線である。AE-2 と AE-4 は volume あたり 500 KB detector-data 制約下で
最も実装しやすく、多数 volume へ複製した場合にも総コストを抑えやすい候補である。
AE-1 は flatten 系の容量不足を切り分ける候補であり、AE-3a は recurrent operator が
MNN 上で許容できるかを測る候補である。AE-3 と AE-5 では、GRU は時系列特徴を
各 frame に付与するために使い、AE-4 と AE-5 では Conv1D は局所時間特徴を
多チャネルに展開して抽出するために使う。いずれも、情報を減らす処理は
TimeDistributed Dense bottleneck に明示的に任せる。
AE-5 は元の CNN-GRU 仮説の検証用として残すが、単純候補が十分に弱いことを
確認する前に既定解にしない。

表 \ref{tab:model-details} は、初期評価対象を TensorFlow/Keras の
\texttt{model.summary()} に近い形式で固定したものである。TD は
\texttt{TimeDistributed}、出力 shape の \texttt{None} は
\texttt{B\_batch} を表す。GRU の parameter count は TensorFlow/Keras の
default である \texttt{reset\_after=True} を仮定する。

\small
\setlength{\tabcolsep}{3pt}
\begin{longtable}{@{}p{0.08\linewidth}p{0.45\linewidth}p{0.24\linewidth}r@{}}
  \caption{主要候補の TensorFlow/Keras 形式 summary}
  \label{tab:model-details} \\
  \toprule
  ID & Layer (Keras syntax) & Output shape & Param \# \\
  \midrule
  \endfirsthead
  \caption[]{主要候補の TensorFlow/Keras 形式 summary（続き）} \\
  \toprule
  ID & Layer (Keras syntax) & Output shape & Param \# \\
  \midrule
  \endhead
  \midrule
  \multicolumn{4}{r}{次ページへ続く} \\
  \endfoot
  \bottomrule
  \endlastfoot
  AE-0 & \texttt{Input(shape=(12, 12))} & \texttt{(None, 12, 12)} & 0 \\
  AE-0 & \texttt{Flatten()} & \texttt{(None, 144)} & 0 \\
  AE-0 & \texttt{Dense(32, activation="relu")} & \texttt{(None, 32)} & 4,640 \\
  AE-0 & \texttt{Dense(8, activation="relu", name="z")} & \texttt{(None, 8)} & 264 \\
  AE-0 & \texttt{Dense(32, activation="relu")} & \texttt{(None, 32)} & 288 \\
  AE-0 & \texttt{Dense(144, activation="linear")} & \texttt{(None, 144)} & 4,752 \\
  AE-0 & \texttt{Reshape((12, 12))} & \texttt{(None, 12, 12)} & 0 \\
  \multicolumn{3}{r}{AE-0 total} & 9,944 \\
  \midrule
  AE-1 & \texttt{Input(shape=(12, 12))} & \texttt{(None, 12, 12)} & 0 \\
  AE-1 & \texttt{Flatten()} & \texttt{(None, 144)} & 0 \\
  AE-1 & \texttt{Dense(64, activation="relu")} & \texttt{(None, 64)} & 9,280 \\
  AE-1 & \texttt{Dense(16, activation="relu", name="z")} & \texttt{(None, 16)} & 1,040 \\
  AE-1 & \texttt{Dense(64, activation="relu")} & \texttt{(None, 64)} & 1,088 \\
  AE-1 & \texttt{Dense(144, activation="linear")} & \texttt{(None, 144)} & 9,360 \\
  AE-1 & \texttt{Reshape((12, 12))} & \texttt{(None, 12, 12)} & 0 \\
  \multicolumn{3}{r}{AE-1 total} & 20,768 \\
  \midrule
  AE-2 & \texttt{Input(shape=(12, 12))} & \texttt{(None, 12, 12)} & 0 \\
  AE-2 & \texttt{TD(Dense(16, activation="relu"))} & \texttt{(None, 12, 16)} & 208 \\
  AE-2 & \texttt{TD(Dense(8, activation="relu"))} & \texttt{(None, 12, 8)} & 136 \\
  AE-2 & \texttt{Flatten()} & \texttt{(None, 96)} & 0 \\
  AE-2 & \texttt{Dense(24, activation="relu")} & \texttt{(None, 24)} & 2,328 \\
  AE-2 & \texttt{Dense(8, activation="relu", name="z")} & \texttt{(None, 8)} & 200 \\
  AE-2 & \texttt{Dense(24, activation="relu")} & \texttt{(None, 24)} & 216 \\
  AE-2 & \texttt{Dense(96, activation="relu")} & \texttt{(None, 96)} & 2,400 \\
  AE-2 & \texttt{Reshape((12, 8))} & \texttt{(None, 12, 8)} & 0 \\
  AE-2 & \texttt{TD(Dense(16, activation="relu"))} & \texttt{(None, 12, 16)} & 144 \\
  AE-2 & \texttt{TD(Dense(12, activation="linear"))} & \texttt{(None, 12, 12)} & 204 \\
  \multicolumn{3}{r}{AE-2 total} & 5,636 \\
  \midrule
  AE-3a & \texttt{Input(shape=(12, 12))} & \texttt{(None, 12, 12)} & 0 \\
  AE-3a & \texttt{GRU(24, return\_sequences=True)} & \texttt{(None, 12, 24)} & 2,736 \\
  AE-3a & \texttt{TD(Dense(8, activation="relu", name="z\_seq"))} & \texttt{(None, 12, 8)} & 200 \\
  AE-3a & \texttt{TD(Dense(24, activation="relu"))} & \texttt{(None, 12, 24)} & 216 \\
  AE-3a & \texttt{GRU(24, return\_sequences=True)} & \texttt{(None, 12, 24)} & 3,600 \\
  AE-3a & \texttt{TD(Dense(12, activation="linear"))} & \texttt{(None, 12, 12)} & 300 \\
  \multicolumn{3}{r}{AE-3a total} & 7,052 \\
  \midrule
  AE-4 & \texttt{Input(shape=(12, 12))} & \texttt{(None, 12, 12)} & 0 \\
  AE-4 & \texttt{Conv1D(24, 3, padding="same", activation="relu")} & \texttt{(None, 12, 24)} & 888 \\
  AE-4 & \texttt{Conv1D(24, 3, padding="same", activation="relu")} & \texttt{(None, 12, 24)} & 1,752 \\
  AE-4 & \texttt{TD(Dense(8, activation="relu", name="z\_seq"))} & \texttt{(None, 12, 8)} & 200 \\
  AE-4 & \texttt{TD(Dense(24, activation="relu"))} & \texttt{(None, 12, 24)} & 216 \\
  AE-4 & \texttt{Conv1D(24, 3, padding="same", activation="relu")} & \texttt{(None, 12, 24)} & 1,752 \\
  AE-4 & \texttt{Conv1D(12, 3, padding="same", activation="linear")} & \texttt{(None, 12, 12)} & 876 \\
  \multicolumn{3}{r}{AE-4 total} & 5,684 \\
  \midrule
  AE-5 & \texttt{Input(shape=(12, 12))} & \texttt{(None, 12, 12)} & 0 \\
  AE-5 & \texttt{Conv1D(24, 3, padding="same", activation="relu")} & \texttt{(None, 12, 24)} & 888 \\
  AE-5 & \texttt{GRU(24, return\_sequences=True)} & \texttt{(None, 12, 24)} & 3,600 \\
  AE-5 & \texttt{TD(Dense(8, activation="relu", name="z\_seq"))} & \texttt{(None, 12, 8)} & 200 \\
  AE-5 & \texttt{TD(Dense(24, activation="relu"))} & \texttt{(None, 12, 24)} & 216 \\
  AE-5 & \texttt{GRU(24, return\_sequences=True)} & \texttt{(None, 12, 24)} & 3,600 \\
  AE-5 & \texttt{TD(Dense(12, activation="linear"))} & \texttt{(None, 12, 12)} & 300 \\
  \multicolumn{3}{r}{AE-5 total} & 8,804 \\
\end{longtable}
\setlength{\tabcolsep}{6pt}
\normalsize

AE-5 は volume あたり 500 KB detector-data 制約内で CNN-GRU 仮説を試す最小構成として、以下の制約を置く。
feature layout は scalar-only の $D=12$ に固定し、CNN は feature 方向ではなく
時間方向の \texttt{Conv1D(24, 3)} として使う。これは 10 秒フレーム列の局所的な
増減を複数チャネルへ展開して抽出するためであり、平均値しかない入力に存在しない
bucket locality を仮定しない。GRU は encoder/decoder とも 1 層、hidden size 24 に固定するが、
系列全体を 1 vector に潰すためには使わない。GRU は各 frame へ文脈を付与し、
情報圧縮は \texttt{TD(Dense(8))} の時系列 bottleneck が担う。Conv1D は特徴抽出、
Dense は圧縮という役割を分ける。Conv2D、
transposed convolution、多層 GRU、attention、モデル内 thresholding は使わない。
この条件では weight は約 36 KB に収まるが、入力統計/state を含む per-volume
500 KB detector-data 達成と aggregate device-fit は変換後の実測で判定する。

すべての候補は、入力と同じ $[1,N,D]$ の正規化統計列を線形出力として再構成する。
alert 判定、しきい値処理、channel-wise error 集計は MNN graph の外で実行する。
これにより、モデルを小さく保ちつつ、運用上のしきい値や説明情報を後から変更できる。
初期実装では分布用 softmax や channel ごとの特殊 head を入れず、正規化済み scalar を
直接復元する。これにより、入力にない histogram 情報を暗黙に仮定しない。

図 \ref{fig:pipeline} に、最終的な検出 pipeline の概念図を示す。

\begin{figure}[tb]
  \centering
  \resizebox{\linewidth}{!}{%
  \begin{tikzpicture}[
      node distance=9mm,
      box/.style={draw, rounded corners=2pt, align=center, minimum height=9mm, minimum width=28mm},
      small/.style={draw, rounded corners=2pt, align=center, minimum height=8mm, minimum width=24mm},
      arrow/.style={-Latex, thick}
    ]
    \node[box] (io) {SCSI/NVMe\\I/O metadata};
    \node[box, right=of io] (stats) {10 秒\\統計化};
    \node[box, right=of stats] (seq) {$N$ frames\\$N \times 10$ s};
    \node[box, right=of seq] (ae) {tiny AE\\MNN CPU};
    \node[box, right=of ae] (score) {再構成誤差\\channel breakdown};
    \node[small, below=of ae] (cal) {benign\\calibration};
    \node[small, below=of score] (alert) {しきい値\\alert};

    \draw[arrow] (io) -- (stats);
    \draw[arrow] (stats) -- (seq);
    \draw[arrow] (seq) -- (ae);
    \draw[arrow] (ae) -- (score);
    \draw[arrow] (cal) -- (score);
    \draw[arrow] (score) -- (alert);
  \end{tikzpicture}
  }
  \caption{ストレージ装置内 AE 検出 pipeline}
  \label{fig:pipeline}
\end{figure}

\subsection{volume あたり 500 KB detector-data 予算}

本研究では、約 2000 volume に約 1 GB の検出器データ領域を割り当てる想定から、
volume あたり約 500 KB を持続的な detector-data 予算として置く。この予算には
model weight 情報と、入力 10 秒統計列、normalization constant、threshold、
score history などの入力統計/state を含める。MNN runtime、共有 library、共通 heap は
多数 volume で償却されるため除外する。activation tensor、operator workspace、
output tensor、inference slot は transient scratch として別に測定する。
初期の目標配分を表 \ref{tab:memory-budget} に示す。

\begin{table}[tb]
  \centering
  \caption{detector-data の初期メモリ目標}
  \label{tab:memory-budget}
  \begin{tabular}{lr}
    \toprule
    要素 & 目標 \\
    \midrule
    model weights and quantization metadata & $\leq 360\,\mathrm{KB}$ \\
    retained 10-second input statistics & $\leq 64\,\mathrm{KB}$ \\
    normalization, threshold, calibration state & $\leq 24\,\mathrm{KB}$ \\
    score history and alert explanation state & $\leq 16\,\mathrm{KB}$ \\
    volume metadata and implementation margin & $\leq 36\,\mathrm{KB}$ \\
    \midrule
    合計 & $\leq 500\,\mathrm{KB}$ \\
    \bottomrule
  \end{tabular}
\end{table}

この配分は、volume あたりの detector data を抑えるための仮配分であり、
最終判断は MNN 変換後の weight 表現と入力統計/state の実測に基づく。
ストレージ装置で $V$ 個の volume/namespace に同時適用する場合、weight を共有できる
実装でも、入力 sequence buffer、normalization constant、threshold、score history は
volume ごとに複製され得る。したがって実験では、単一 volume の detector-data fit だけでなく、
$V \simeq 2000$ を含む aggregate memory と 10 秒あたり CPU 時間を報告する。
特に $N$ は、retained input statistics と recurrent/temporal operator の
transient scratch に直接影響するため、検出性能だけでなく volume あたり 500 KB detector-data 予算との
trade-off として選ぶ。
MNN は on-device/embedded inference、CPU backend、ONNX などからの変換、
MNN-Compress、FP16/Int8 量子化を提供する\cite{mnn}。ただし、runtime の軽量性は
detector-data fit の保証ではないため、変換、量子化、parity、per-volume data、
transient scratch の
測定を独立した評価項目にする。
